\chapter{Documentation}\label{ch:documentation}
This chapter documents the hard- and software of the RCSL, wherein the hardware section describes the hardware components, which the device is built from and the software section explains the architecture of the software and its functionality.

\section{Hardware}
The figure below shows the hardware components, used in the project, are a RaspberryPI (RPI), a USB network card, an ESP32 microcontroller, a display and a rotary encoder.

\begin{figure}[h]
    \centering
    \includegraphics[width=0.8\linewidth]{figures/Abbildungen/Devices.png}
    \caption{Hardware Components}
    \label{fig:components}
\end{figure}

As the main computer the RaspberryPI 4B with 8GB of RAM was chosen due to its ease of use, linux based operating system and widely available documentation. 

The RPI has an integrated WiFi and Bluetooth card, as well as onboard ethernet networking.
However to support all WiFi standards an external network card is needed, since the internal card of the RPI does not support the ciphers used by WPA3.
The Fenvi AX1800 USB WiFi card was chosen for its compatibility with the linux kernel, according to \cite{morrownr_USB_WiFi}, and low cost.

For some types of attack, such as KRACK\footnote{see \cref{ch:testing}}, a WiFi station, which is connected to the test network is needed, which in this case is an ESP32 microcontroller unit (MCU).
Selecting from the large offer of ESP32 MCUs, again size was taken into account and the ESP32-C6 Mini by Waveshare was chosen for its small size and good feature set.
It connects to the RPI's GPIO pins using simple jumper cables, which supply the ESP32 with power and establish an UART connection.

Another device connected to the RPI's GPIO header is the rotary encoder.
On its five pins, it receives 3.3V supply voltage and outputs a gray code, indicating the direction of rotation, on two of its outputs. 
The fifth pin is used as the input for the pushbutton action of the encoder.

The hardware components are mounted in a 3d printed case, like shown in \cref{fig:RCSL}
%\begin{figure}[h]
%    \centering
%    \includegraphics[width=\linewidth]{figures/Abbildungen/Devices.png}
%    \caption{RCSL within its case}
%    \label{fig:RCSL}
%\end{figure}

\section{Software}

\textit{Note: Due to my limited experience in software development, the project does not focus on performance or software quality.
The software is written to implement the wanted functionality and should be seen as experimental, since it is neither optimised for performance, nor have proper error handling mechanisms been implemented.}

\subsection{Overview}
The software of the project is mostly modular and it consists of a central program, which runs the user interface (UI) and then executes various shell scripts, with which the functionality is implemented.
As to be seen in \cref{fig:architecture}, besides the UI program, there are programs for MQTT communication and input processing, also the services of the Network Manager and Juice Shop project are used. 

\begin{figure}[h]
    \includegraphics[width=\linewidth]{figures/Abbildungen/architecture.pdf}
    \centering
    \caption{Software Architecture}
    \label{fig:architecture}
\end{figure}

\textit{Note: only the UI program is written object oriented and only Menu and MenuOption are object classes.}


\subsection{UI Program}

Upon booting the device, a systemd service executes start.sh, which starts the encoder program and the UI program with the output of the encoder program piped into the input of the UI.
The UI displays the main menu, depicted in \cref{fig:menu_main}, which the user can interact with by using the rotary encoder.
Each page of the UI is a "Menu" object, each menu has a title, options stored inside a vector of the type "MenuOption", a variable to track the currently selected option, a method to add options to a menu and methods for navigation, input capture and displaying.
%The program starts by calling the display() function of the main menu, which then prints the options of the menu periodically, highlighting the selected option in orange.
MenuOption objects have a name, as well as an action of the std::function type, which can be executed by calling the execute() method.
In the main of the program, the constructors for menues and the addOption() methods are beeing used, to create the different menues and their submenues.
An overview of the menues and options, currently present inside the UI program can be seen in \cref{tab:menues_options}.
The program starts by calling the navigate() function of the main menu, which then displays all available submenues and menu options.
All submenues work in the same way as the main menu, getting their respective navigate() called upon selection.
During navigation the UI peridiodically prints the available options, using unicode to highlight the currently selected option as well as to give contrast to the title of the menu.

\begin{figure}[h]
    \includegraphics[width=0.3\linewidth]{figures/Abbildungen/hackerypi/menu_main.png}
    \centering
    \caption{Main Menu}
    \label{fig:menu_main}
\end{figure}

\textit{More examples for the UI can be found in the appendix, for example \cref{fig:menues}.}
    
\subsection{Encoder Program}
For input the device uses a rotary encoder with a pushbutton, turning the encoder moves the menu selection up or down, the pushbutton executes the selected option.
The encoder is connected to the RaspberryPIs GPIO interface, its rotation is segmented with clicks, similar to the scroll wheel of a computer mouse
Rotation of the encoder is segemented with physical clicks outputs a grey code, on the signals DT and CLK (see \cref{fig:encoder_processing}).
Like shown, rotation between two clicks outputs one complete grey code sequence, which is processed by the encoder program.
This happens within the rotation handling interrupt service routine (ISR), shown in \cref{lst:ISR_rotation}, which is executed, when DT outputs a high-low transition.

\lstinputlisting[
    firstline=13,
    lastline=24,
    language=C,
    caption={Rotation handling ISR}
    label={lst:ISR_rotation}
]
    {listings/Raspi-CyberSec-Lab-Project/Program/encoder.h}

The ISR checks the state of CLK, which determines the direction, the encoder was turned: CLK=0 means, the encoder has been turned clockwise, CLK=1 means counter-clockwise (see \cref{fig:encoder_processing}).
Subsequently a command in form of a string, representing an ANSI Escape Code, is printed.
Escape Codes are used as an output for peripheral buttons, which are not a character, for example \lstinline[]|"\033[A"| describes the pressing of the downwards button on a keyboard.


For example, the User turns the the encoder one click clockwise, therefore main receives one downwards arrow symbol and the selections moves to the next option below.
Because the commands used for navigation are ordinary ANSI escape codes, the program can be easily tested on any computer by navigating with the keyboard.

\subsection{Shell Scripts}
When an action is executed, main calls the according shell script with the system() function.
The wifiActivate.sh script is used to open wifi hotspots, where the user can choose the type of Wifi standard beeing used. 
It receives an input command and then establishes a WEP, WPA, WPA2 or WPA3 network.

\lstinputlisting[
    language=bash,
    caption={extract of wifiActivate.sh},
    lastline=16,
    breaklines=true,
    basicstyle=\ttfamily\fontsize{8}{9}\selectfont
]
    {listings/Raspi-CyberSec-Lab-Project/Skripte/wifiActivate.sh}
    
With the wifiReset.sh script, the user can deactivate the active network.
Once a network has been cracked, the wifiNewPassword.sh script can be used to give a new random password to the selected network.
The script the pulls a random password out of a password list, which contains popular passwords that have been filtered to work with the specification of each network standard.
To display the current status of the network, meaning its security standard, SSID and password, the wifiStatus.sh script is used.


\begin{figure}[h]
    \centering
    \includegraphics[width=0.24\linewidth]{figures/Abbildungen/hackerypi/status_wifi.png}
    \caption{WiFi status menu}
    \label{fig:menu_status}
\end{figure}

After a WiFi hotspot has been activated, a command with the connection details is sent to the ESP32 over uart.
This functionality is implemented by the uartSendNI.sh script, which gets called from wifiActivate.sh.
Receiving the network information, the ESP32 will then connect to the provided network and act as a communication partner for the RPI.

Selecting the "monitor" option, the user is presented with the options of de-/activating the network monitoring as well as viewing or deleting the monitoring log.
When activating monitoring, wifiMonitor.sh uses the \lstinline[]|iw event -T| command to log devices connecting and disconnecting from the network.

\lstinputlisting[
    language=bash,
    caption={extract of wifiActivate.sh},
    breaklines=true,
    basicstyle=\ttfamily\fontsize{8}{9}\selectfont
]
    {listings/Raspi-CyberSec-Lab-Project/Skripte/uartSendNI.sh}

Inside the webapp menu the JuiceShop and MQTT communication can be found.
When activating the JuiceShop, the juiceShop.sh script will start up the server for the OWASP JuiceShop project and open a hotspot, which can be used to access the simulated webstore.

When a WiFi network has already been activated from the WiFi menu, a simulated MQTT communication can be started from the MQTT menu.
Starting a conversation will again send a command via uart to the ESP32, prompting it to connect to the MQTT broker, running on the RPI.
Following up uart.sh starts the uart program, which starts an infinite MQTT conversation with the ESP32.

The development menu contains options used during development, these can be adjusted on the fly in the development.sh script.

When selecting the shutdown option, all running processes are stopped and the shutdown.sh script calls \lstinline[]|shutdown now| to turn off the device.




